%!TEX TS-program = xelatex
%!TEX encoding = UTF-8 Unicode
% Generated by imessage-book. Build with `tectonic book.tex` or `xelatex book.tex`.
% Requires xelatex/tectonic for Unicode + emoji font support.

\documentclass[10pt,oneside]{memoir}
\setstocksize{<< page.height >>}{<< page.width >>}
\settrimmedsize{\stockheight}{\stockwidth}{*}
\setlrmarginsandblock{<< page.margin >>}{<< page.margin >>}{*}
\setulmarginsandblock{<< page.margin >>}{*}{*}
\checkandfixthelayout
% Chat content is a stack of rigid, unbreakable bubbles with little stretchable glue, so
% let page bottoms fall where content ends rather than stretching gaps to justify them.
\raggedbottom

\usepackage{fontspec}
\usepackage{microtype}
\usepackage{graphicx}
\usepackage[export]{adjustbox}
\usepackage{xcolor}
\usepackage{tikz}
\usepackage[most]{tcolorbox}
\usepackage{varwidth}

% Bubble colors (overridable via book.toml [theme]).
\definecolor{msgblue}{HTML}{<< theme.me_color >>}
\definecolor{msggray}{HTML}{<< theme.them_color >>}
\definecolor{msggreen}{HTML}{<< theme.sms_color >>}
\definecolor{metagray}{HTML}{<< theme.meta_color >>}

% Emoji font: prefer a configured font (e.g. Apple Color Emoji, rendered via HarfBuzz),
% then Noto Emoji, and finally degrade to a no-op so the document always builds.
<% if emoji_font %>\IfFontExistsTF{<< emoji_font >>}
  {\newfontfamily\emojifont{<< emoji_font >>}[Renderer=Harfbuzz]}
  {\IfFontExistsTF{Noto Emoji}{\newfontfamily\emojifont{Noto Emoji}}{\newcommand{\emojifont}{}}}
<% else %>\IfFontExistsTF{Noto Emoji}
  {\newfontfamily\emojifont{Noto Emoji}}
  {\newcommand{\emojifont}{}}
<% endif %>

<% if theme.font %>\IfFontExistsTF{<< theme.font >>}{\setmainfont{<< theme.font >>}}{}<% endif %>

% Running header shows the current month/chapter.
\makepagestyle{messagebook}
\makeevenhead{messagebook}{\thepage}{}{\slshape\leftmark}
\makeoddhead{messagebook}{\slshape\rightmark}{}{\thepage}
\pagestyle{messagebook}

% Bubbles hug their content up to this width; photos are capped a little narrower so a
% single tall image never dominates the page.
\newlength{\bubblemax}\setlength{\bubblemax}{0.62\linewidth}
\newlength{\photomax}\setlength{\photomax}{0.5\linewidth}

\tcbset{
  bubblebase/.style={boxrule=0pt,arc=8pt,left=10pt,right=10pt,top=5pt,bottom=5pt,
    before skip=0pt,after skip=0pt,nobeforeafter},
  mestyle/.style={bubblebase,colback=msgblue,coltext=white,colframe=msgblue},
  themstyle/.style={bubblebase,colback=msggray,coltext=black,colframe=msggray},
  smsstyle/.style={bubblebase,colback=msggreen,coltext=white,colframe=msggreen},
  % Reaction/tapback pill: a small white bubble like the one Messages floats on a bubble.
  pillstyle/.style={boxrule=0.4pt,arc=6pt,left=5pt,right=5pt,top=1pt,bottom=1pt,
    before skip=0pt,after skip=0pt,nobeforeafter,colback=white,coltext=black,colframe=msggray},
}

% A colored text bubble, capped at \bubblemax and hugging its content. #1 = style.
\newcommand{\textbubble}[2]{\tcbox[#1]{\begin{varwidth}{\bubblemax}#2\end{varwidth}}}

% A photo that *is* the bubble: rounded corners, no colored frame, no caption.
% #1 = extra tikz overlay code (e.g. a play button), #2 = image path.
\newsavebox{\mbphoto}\newlength{\mbphw}\newlength{\mbphh}
\newcommand{\photobox}[2]{%
  \sbox{\mbphoto}{\includegraphics[width=\photomax,height=0.32\textheight,keepaspectratio]{#2}}%
  \settowidth{\mbphw}{\usebox{\mbphoto}}%
  \settoheight{\mbphh}{\usebox{\mbphoto}}%
  \begin{tikzpicture}
    \clip[rounded corners=12pt] (0,0) rectangle (\mbphw,\mbphh);
    \node[anchor=south west,inner sep=0pt] at (0,0){\usebox{\mbphoto}};
    #1
  \end{tikzpicture}%
}
\newcommand{\photobubble}[1]{\photobox{}{#1}}
% Same, with a translucent play button centered — for video poster frames.
\newcommand{\videobubble}[1]{\photobox{%
  \coordinate (mbc) at (0.5\mbphw,0.5\mbphh);
  \fill[black,opacity=0.40] (mbc) circle[radius=10pt];
  \fill[white] ([xshift=-3.5pt,yshift=6pt]mbc) -- ([xshift=-3.5pt,yshift=-6pt]mbc) -- ([xshift=6pt]mbc) -- cycle;
}{#1}}

% Right/left aligned message rows (me / them). #1 = "r" or "l", #2 = box content.
%
% Page-break control keeps each message intact without welding a whole back-and-forth
% into one unbreakable block. A welded burst that runs near a full page tall strands the
% gap marker / day / chapter heading before it on an otherwise-empty page, and tall photo
% runs can't share a page — the cause of the alternating full/blank pages. So \brow
% forbids a break *before* the row (a bubble stays glued to a preceding sender name or
% reply quote), while \pbrow permits one and is used for photos/videos, letting a tall
% image start a fresh page instead of shoving the whole run down.
\newcommand{\brow}[2]{\par\nobreak\vspace{1.5pt}\noindent\makebox[\linewidth][#1]{#2}\par}
\newcommand{\pbrow}[2]{\par\penalty0\vspace{1.5pt}\noindent\makebox[\linewidth][#1]{#2}\par}
% Small metadata line (timestamp / reactions) hugging the same edge as the bubble; the
% leading \nobreak keeps it on the same page as the bubble it annotates.
\newcommand{\mrow}[2]{\par\nobreak\noindent\makebox[\linewidth][#1]{\scriptsize\color{metagray}#2}\par}
% Permit a page break between two messages: individual messages stay whole, but a long
% exchange flows across pages instead of forcing the whole run onto one.
\newcommand{\msgsep}{\par\penalty0}

\newcommand{\attachlabel}[1]{{\footnotesize #1}}
\newcommand{\replyquote}[1]{{\footnotesize\itshape\color{metagray}#1}}
\newcommand{\apptag}[1]{{\footnotesize\textbf{#1}\par}}
\newcommand{\tapbackpill}[1]{\tcbox[pillstyle]{\scriptsize #1}\hspace{2pt}}
\newcommand{\sendername}[1]{\par\addvspace{5pt}\noindent{\scriptsize\color{metagray}#1}\par\nobreak}
\newcommand{\dayheading}[1]{\begin{center}\vspace{8pt}{\footnotesize\itshape\color{metagray}#1}\vspace{2pt}\end{center}}
\newcommand{\gapmark}[1]{\par\addvspace{5pt}\begin{center}{\scriptsize\color{metagray}#1}\end{center}\addvspace{3pt}}
\newcommand{\systemline}[1]{\par\addvspace{5pt}\begin{center}{\scriptsize\itshape\color{metagray}#1}\end{center}\addvspace{5pt}}

<% if author %>\author{<< author | tex >>}<% else %>\author{}<% endif %>
\title{<< title | tex >>}
\date{}

%%% BEGIN DOCUMENT
\begin{document}

\frontmatter
\begin{titlingpage}
<% if cover %>\begin{center}\includegraphics[max width=0.8\textwidth,max height=0.5\textheight]{<< cover >>}\par\bigskip\end{center}<% endif %>
\maketitle
\end{titlingpage}

<% if dedication %>
\movetooddpage
\begin{vplace}
\begin{center}\textit{<< dedication | tex >>}\end{center}
\end{vplace}
<% endif %>

<% if stats %>
\movetooddpage
\begin{center}{\Large\bfseries By the Numbers}\end{center}
\bigskip
\begin{description}
\item[Messages] << stats.total >> (<< stats.from_me >> from you, << stats.from_others >> received)
\item[Span] << stats.first_date | tex >> to << stats.last_date | tex >> (<< stats.days >> day<% if stats.days != 1 %>s<% endif %>)
\item[Words sent] << stats.words >>
\item[Photos] << stats.photos >>
<% if stats.busiest_day %>\item[Busiest day] << stats.busiest_day | tex >>
<% endif %><% if stats.busiest_hour %>\item[Busiest hour] << stats.busiest_hour | tex >>
<% endif %><% if stats.longest_streak > 1 %>\item[Longest streak] << stats.longest_streak >> days
<% endif %><% if stats.top_emoji %>\item[Top emoji] <% for e in stats.top_emoji %>{\emojifont << e.emoji >>}~<< e.count >>\quad <% endfor %>
<% endif %><% if is_group %>\item[Most active] <% for s in stats.per_sender %><< s.name | tex >>: << s.count >>\newline <% endfor %>
<% endif %>\end{description}
<% endif %>

\movetooddpage
\tableofcontents

\cleartorecto
\mainmatter
\chapterstyle{dash}
\addtolength\afterchapskip{1cm}

<% for chapter in chapters %>
\chapter{<< chapter.title >>}
<% for day in chapter.days %>
\dayheading{<< day.heading | tex >>}
<% for m in day.messages %>
<% if m.system %>\systemline{<< m.system | tex >>}
<% else %>\msgsep
<% if m.gap_before %>\gapmark{<< m.gap_before | tex >>}
<% endif %>
<% if m.sender %>\sendername{<< m.sender | tex >>}
<% endif %>
<% set bstyle = "smsstyle" if (m.from_me and m.service == "sms") else ("mestyle" if m.from_me else "themstyle") %>
<% set al = "r" if m.from_me else "l" %>
<% if m.reply_to %>\brow{<< al >>}{\replyquote{<< m.reply_to | tex >>}}
<% elif m.reply %>\brow{<< al >>}{\replyquote{Reply}}
<% endif %>
<% for a in m.attachments %>
<% if a.src and a.kind == "video" %>\pbrow{<< al >>}{\videobubble{<< a.src >>}}
<% elif a.src and a.kind == "image" %>\pbrow{<< al >>}{\photobubble{<< a.src >>}}
<% elif a.kind == "image" or a.kind == "video" %><# a photo/video we couldn't embed: skip it rather than print a bare "Photo"/"Video" label #>
<% else %>\brow{<< al >>}{\textbubble{<< bstyle >>}{\attachlabel{<< a.label | tex >><% if a.caption %> \textperiodcentered\ << a.caption | tex >><% endif %>}}}
<% endif %>
<% endfor %>
<% if m.text or m.app %>\brow{<< al >>}{\textbubble{<< bstyle >>}{<% if m.app %>\apptag{<< m.app | tex >>}<% endif %><% if m.text %><< m.text | tex >><% endif %>}}
<% endif %>
<% if m.tapbacks %>\brow{<< al >>}{<% for t in m.tapbacks %>\tapbackpill{<< t | tex >>}<% endfor %>}
<% endif %>
\mrow{<< al >>}{<< m.time >><% if m.edited %> \textperiodcentered\ edited<% endif %><% if m.effect %> \textperiodcentered\ << m.effect | tex >><% endif %>}
<% endif %>
<% endfor %>
<% endfor %>
<% endfor %>

<% if gallery %>
\appendix
\chapter{Photos}
\noindent
<% for src in gallery %>\includegraphics[max width=0.3\textwidth,max height=0.25\textheight]{<< src >>}\hspace{3pt}<% endfor %>
<% endif %>

\end{document}
